\documentclass{article}
\usepackage{amsmath}
\usepackage{amsthm}

\newtheorem{theorem}{Problem}

\begin{document}

\begin{theorem}
Find all solutions to\n\[\sin \left( \tan^{-1} (x) + \cot^{-1} \left( \frac{1}{x} \right) \right) = \frac{1}{3}.\]Enter all the solutions, separated by commas.
\end{theorem}

\begin{proof}
Step 0: Let's set $y = \tan^{-1} x$.
Step 1: Then $x = \tan y$, and $\frac{1}{x} = \cot y$.
Step 2: That's true. So the expression becomes $\sin (y + \cot^{-1} (\cot y))$.
Step 3: Since $\cot^{-1} (\cot y) = y$, the expression simplifies to $\sin (y + y)$.
Step 4: So, it becomes $\sin (2y)$.
Step 5: That's correct. And we want to find values of $y$ that satisfy the equation $\sin (2y) = \frac{1}{3}$.
Step 6: Yes, let's find the solutions to that equation.
Step 7: The solutions are $2y = \arcsin \frac{1}{3} + 2k\pi$ and $2y = \pi - \arcsin \frac{1}{3} + 2k\pi$, where $k$ is an integer.
Step 8: Those are the solutions for $y$. But we want to find the solutions for $x$.
Step 9: To do that, we just need to use the fact that $x = \tan y$.
Step 10: Right. So let's plug in the solutions for $y$ into $x = \tan y$.
Step 11: The solution $2y = \arcsin \frac{1}{3} + 2k\pi$ gives $x = \tan (\frac{1}{2} \arcsin \frac{1}{3} + k\pi)$.
Step 12: And the other solution $2y = \pi - \arcsin \frac{1}{3} + 2k\pi$ gives $x = \tan (\frac{1}{2} (\pi - \arcsin \frac{1}{3}) + k\pi)$.
Step 13: Which simplifies to $x = \tan (\frac{\pi}{2} - \frac{1}{2} \arcsin \frac{1}{3} + k\pi)$.
Step 14: Let's simplify those solutions even further.
Step 15: Note that $\tan$ is periodic with a period of $\pi$.
Step 16: Also, $\tan (\frac{\pi}{2} - \theta) = \cot \theta$.
Step 17: So the first solution simplifies to $x = \tan (\frac{1}{2} \arcsin \frac{1}{3})$. What about the second?
Step 18: The second solution simplifies to $x = \cot (\frac{1}{2} \arcsin \frac{1}{3}) = \frac{1}{\tan (\frac{1}{2} \arcsin \frac{1}{3})}$.
Step 19: So the two solutions are $x = \tan (\frac{1}{2} \arcsin \frac{1}{3})$ and $x = \frac{1}{\tan (\frac{1}{2} \arcsin \frac{1}{3})}$.
Step 20: Let's find $\tan (\frac{1}{2} \arcsin \frac{1}{3})$ using the identity $\tan (\frac{1}{2} \theta) = \frac{\sin \theta}{1 + \cos \theta}$.
Step 21: If we plug in $\theta = \arcsin \frac{1}{3}$, we get $\tan (\frac{1}{2} \arcsin \frac{1}{3}) = \frac{\sin (\arcsin \frac{1}{3})}{1 + \cos (\arcsin \frac{1}{3})}$.
Step 22: What are $\sin (\arcsin \frac{1}{3})$ and $\cos (\arcsin \frac{1}{3})$?
Step 23: $\sin (\arcsin \frac{1}{3}) = \frac{1}{3}$. To find $\cos (\arcsin \frac{1}{3})$, we use the identity $\sin^2 \theta + \cos^2 \theta = 1$.
Step 24: If $\theta = \arcsin \frac{1}{3}$, then $\sin \theta = \frac{1}{3}$, so the identity becomes $\frac{1}{9} + \cos^2 \theta = 1$.
Step 25: Which means $\cos^2 \theta = \frac{8}{9}$. Square rooting, we get $\cos (\arcsin \frac{1}{3}) = \pm \frac{2\sqrt{2}}{3}$.
Step 26: But which sign should we choose?
Step 27: We can use the identity $\cos (\arcsin x) = \sqrt{1 - x^2}$ to find the correct sign.
Step 28: If we plug in $x = \frac{1}{3}$, we get $\cos (\arcsin \frac{1}{3}) = \sqrt{1 - \frac{1}{9}} = \sqrt{\frac{8}{9}} = \pm \frac{2\sqrt{2}}{3}$.
Step 29: So the correct sign is $+$. Therefore, $\tan (\frac{1}{2} \arcsin \frac{1}{3}) = \frac{\frac{1}{3}}{1 + \frac{2\sqrt{2}}{3}} = \frac{1}{2\sqrt{2} + 3}$.
Step 30: And $\frac{1}{\tan (\frac{1}{2} \arcsin \frac{1}{3})} = 2\sqrt{2} + 3$.
Step 31: So the two solutions are $x = \frac{1}{2\sqrt{2} + 3}$ and $x = 2\sqrt{2} + 3$.
Step 32: That's the answer.
Step 33: No, the simplification at Step 30 was correct. The value is indeed $\frac{1}{2\sqrt{2} + 3}$. Therefore, the solution gives the boxed value.
\end{proof}

\end{document}
